This thesis has presented two complementary contributions to monocular 3D object detection, each addressing a distinct limitation of prior work.

\textit{MonoDSSMs} tackles the computational inefficiency of Transformer-based detectors by introducing the first Mamba-based architecture for monocular 3D detection. Through the proposed BiMamba2 and CrossMamba modules, together with a multi-scale depth enhancement strategy, the framework captures global context with linear complexity and integrates depth-aware features efficiently. On the KITTI benchmark, MonoDSSMs achieves detection quality comparable to its Transformer-based baseline while using 2.2x fewer parameters and running 1.28x faster. Our plug-and-play study further demonstrates that the framework generalizes well across modern Vision Mamba architectures, accommodating both fixed-scan and learnable-scan strategies.

\textit{Mono3DV} addresses a critical limitation of DETR-like monocular 3D detectors: the exclusion of 3D attributes from bipartite matching, which leads to suboptimal supervision and unstable training. By introducing a 3D-Aware Bipartite Matching strategy, a 3D-DeNoising scheme, and a Variational Query DeNoising mechanism, Mono3DV resolves the misalignment between 2D-only matching criteria and the underlying 3D objective. Without relying on external data, it achieves state-of-the-art results on the KITTI and Waymo benchmarks and provides consistent gains as a plug-and-play module for multi-view detectors on nuScenes.

\section*{Future Directions}

\textbf{MonoDSSMs.} Our experiments on modern Vision Mamba architectures suggest that learnable scanning strategies, such as input-dependent tree topologies and deformable scan offsets, offer a promising path beyond the bidirectional scan adopted in MonoDSSMs. However, current learnable variants come with non-trivial overhead in parameters and inference speed. A natural direction is to design lightweight learnable scanning mechanisms that preserve the input-adaptive behavior responsible for accuracy gains while keeping the efficiency profile that makes MonoDSSMs attractive for real-time deployment. This may involve sparse or low-rank parameterizations of the scan path, hardware-aware scan kernels, or hybrid schemes that combine a fast fixed scan with a small learnable correction.

\textbf{Mono3DV.} The proposed Variational Query DeNoising and Forward-Looking Distillation techniques are generally applicable to any DETR-based model. However, this work has focused on their application to monocular 3D detection. A promising direction is to investigate their broader applicability and generalization to other domains, such as 2D object detection, instance segmentation, and multi-view 3D detection, where similar matching instability and gradient vanishing issues may arise. To further strengthen the denoising process itself, exploring alternative generative models beyond variational autoencoders represents another promising avenue. Normalizing flow-based and diffusion-based denoising approaches, in particular, could offer richer noise distributions and more expressive query refinement.

\section*{Closing Remarks}

Together, the two contributions of this thesis advance monocular 3D object detection along complementary axes of efficient feature modeling and geometry-aware set-based prediction. We hope that the architectural and training innovations presented here will serve as useful building blocks for future research in efficient and reliable 3D perception from monocular cameras.